\documentclass{article}

\usepackage{amsmath,amssymb}
\usepackage{xurl}
\usepackage[hidelinks]{hyperref}
\setlength{\emergencystretch}{2em}

% Shared commands for notes in docs/paper-gaps/.

\DeclareUnicodeCharacter{2080}{\ifmmode{}_{0}\else\textsubscript{0}\fi}
\DeclareUnicodeCharacter{2081}{\ifmmode{}_{1}\else\textsubscript{1}\fi}
\DeclareUnicodeCharacter{2082}{\ifmmode{}_{2}\else\textsubscript{2}\fi}
\DeclareUnicodeCharacter{2099}{\ifmmode{}_{n}\else\textsubscript{n}\fi}

\newcommand{\Lean}{\textsc{Lean}}
\ExplSyntaxOn
\NewDocumentCommand{\leanid}{m}{
  \ifmmode
    \textsf{\detokenize{#1}}
  \else
    \group_begin:
    \urlstyle{sf}
    \tl_set:Nn \l_tmpa_tl {#1}
    \tl_replace_all:Nnn \l_tmpa_tl {\_} {_}
    \exp_args:NV \path \l_tmpa_tl
    \group_end:
  \fi
}
\ExplSyntaxOff
\providecommand{\tr}{\operatorname{tr}}
\newcommand{\tnleanIssueBase}{https://github.com/LionSR/TNLean/issues/}
\newcommand{\tnleanPrBase}{https://github.com/LionSR/TNLean/pull/}
\newcommand{\tnleanDocBase}{https://sirui-lu.com/TNLean/docs/}
\newcommand{\ghissue}[1]{\href{\tnleanIssueBase#1}{GitHub issue \##1}}
\newcommand{\ghpr}[1]{\href{\tnleanPrBase#1}{PR \##1}}
\newcommand{\doclink}[1]{\href{\tnleanDocBase#1.html}{\path{#1}}}

% Dirac notation and the identity operator, as in blueprint/src/macros/common.tex.
% \providecommand keeps note-local definitions authoritative.
\usepackage{dsfont}
\providecommand{\ket}[1]{|#1\rangle}
\providecommand{\bra}[1]{\langle #1|}
\providecommand{\braket}[2]{\langle #1 | #2 \rangle}
\providecommand{\ketbra}[2]{|#1\rangle\!\langle #2|}
\providecommand{\Id}{\mathds{1}}
\providecommand{\one}{\mathds{1}}
\providecommand{\C}{\mathbb{C}}
\providecommand{\MN}[1]{M_{#1}(\C)}
\providecommand{\MD}{M_D(\C)}

% External referencing for paper sources.  The build helper writes sanitized
% auxiliary files under build/paper-aux/ and excludes duplicated source labels.
\usepackage{xr}

% Import a generated paper auxiliary file with a caller-supplied label prefix.
% The candidate paths support builds from docs/paper-gaps/, the repository root,
% and build/paper-aux/ itself.
\newcommand{\importpaperaux}[2]{%
  \IfFileExists{../../build/paper-aux/#2.aux}{%
    \externaldocument[#1]{../../build/paper-aux/#2}%
  }{%
    \IfFileExists{build/paper-aux/#2.aux}{%
      \externaldocument[#1]{build/paper-aux/#2}%
    }{%
      \IfFileExists{#2.aux}{%
        \externaldocument[#1]{#2}%
      }{}%
    }%
  }%
}

% General external reference macros
\newcommand{\paperref}[2]{\ref{#1-#2}}
\newcommand{\papereqref}[2]{(\ref{#1-#2})}

% CPSV16 source (1606.00608 / MPDO-22-12-17-2.tex) external document and shortcuts
\importpaperaux{cpsv16-}{1606.00608/MPDO-22-12-17-2-tnlean}
\newcommand{\cpsvref}[1]{\paperref{cpsv16}{#1}}
\newcommand{\cpsveqref}[1]{\papereqref{cpsv16}{#1}}

% Verdict marker: \gapnote{<kind>}{<status>}. Kind is the policy
% classification (clarification, local-correction, scope-restriction,
% unfaithful, false-source, open-gap); status is open, wip (elimination
% underway), resolved, or historical. Machine-read by the site index;
% prints a verdict line.
\newcommand{\gapnote}[2]{\par\noindent\textbf{Verdict:} #1 (#2).\par}


\title{Historical Note: Non-Dominant Per-Block Projection in CPSV16 \S II Step~1}
\author{}
\date{2026-05-13}

\begin{document}
\maketitle
\gapnote{unfaithful}{historical}

\section{Status}

\textbf{Historical proof archaeology and a forbidden route, not an active gap.}
The corrected proportional fundamental theorem is
\leanid{MPSTensor.fundamentalTheorem_proportional_canonicalForm}.  It compares
active BNT families whose listed copy weights are nonzero and obtains a
bijection of the basis sectors with gauge-phase equivalences.  The unrestricted
statement printed in CPSV16 remains false: its weak BNT convention allows an
unrelated normal representative whose coefficient is identically zero, so the
claimed equality of BNT cardinalities fails.  The analysis below is retained
because it explains why non-dominant per-block projection cannot prove the
proportional theorem and must not be revived.

\section{Notation}

Fix two families of MPS site tensors $\{A_j\}_{j=1}^{r_A}$ and
$\{B_k\}_{k=1}^{r_B}$ with nonzero scalars
$\mu^A_j,\mu^B_k\in\mathbb C$. The assembled MPV at length~$N$ is
\begin{align}
    V^{(N)}(A_{\mathrm{tot}})
       &=\sum_{j=1}^{r_A}(\mu^A_j)^N\,V^{(N)}(A_j),
    & V^{(N)}(B_{\mathrm{tot}})
       &=\sum_{k=1}^{r_B}(\mu^B_k)^N\,V^{(N)}(B_k),
\end{align}
written $V^{(N)}(\mathtt{toTensorFromBlocks} \mu A)$ in the
formalization.\footnote{%
    \leanid{MPSTensor.toTensorFromBlocks};
    \leanid{MPSTensor.mpv_toTensorFromBlocks_eq_sum}.}
A cross-overlap $\langle V^{(N)}(A_j),V^{(N)}(B_k)\rangle$ is
\leanid{MPSTensor.mpvOverlap}. \emph{Eventual proportionality} of the
assembled families is the hypothesis that for some scalar sequence
$c_N\ne 0$ cofinally,
\begin{align}
    \sum_{j}(\mu^A_j)^N V^{(N)}(A_j)
    &= c_N\sum_{k}(\mu^B_k)^N V^{(N)}(B_k)
    \qquad(\forall^{\mathrm{cof}}N),
    \label{eq:prop}
\end{align}
formalized as
\leanid{EventuallyNonzeroProportionalMPV}\textsubscript{2}.

The retired one-copy normal-canonical BNT route carried a single weight
$\mu_k$ per sector with the strict ordering
\begin{align}
    |\mu_1|>|\mu_2|>\cdots>|\mu_r|.
    \tag{\texttt{mu\_strict\_anti}}
\end{align}
Together with the dominant normalization $|\mu_1|=1$, every sub-dominant
block had modulus strictly less than~$1$.  These fields are not part of the
current \leanid{IsNormalCanonicalFormBNT} predicate, which allows equal-modulus
blocks.  The discussion below concerns the deleted strict-modulus variant.
\cite[Definition~4.2 and lines~1864--1884]{Cirac2021Matrix} states the source
structure using repeated copies and raw weights $\mu_{j,q}$ with sector
multiplicities $r_j\ge 1$.  The current \leanid{SectorDecomposition} surface
retains those power-sum coefficients; see
\path{docs/paper-gaps/cpsv16_ft_one_copy_scope_restriction.tex}.

\section{Cited assertion}

The block-selection paragraph in the proof of Theorem~\texttt{thm1}
\cite[line~1182]{Cirac2016MPDO_arXiv}
fixes any $B$-block $B_{k_0}$ and rules out
\begin{align}
    \langle V^{(N)}(A_j),V^{(N)}(B_{k_0})\rangle
    \xrightarrow[N\to\infty]{}0
    \qquad(\text{all }j),
    \tag{$\mathcal H_{k_0}$}
    \label{eq:Hk0}
\end{align}
by Corollary~\texttt{Lem1} \cite[line~1131]{Cirac2016MPDO_arXiv}, which
gives eventual linear independence of
$\{V^{(N)}(A_j)\}_j\cup\{V^{(N)}(B_{k_0})\}$ from~\eqref{eq:Hk0}.  The
source then treats this as contradicting the proportionality identity at
\cite[eq.~\texttt{eq:II\_ABasicTensors}, line~286]{Cirac2016MPDO_arXiv}.
This is the fixed-block claim made at line~1182.  The passage does not assume
all-pairwise cross-family decay and does not establish linear independence of
the full combined family.

\section{Retired formal target}

The historical target was, for $k_0:\mathrm{Fin}\,r_B$ arbitrary: given the
one-copy normal-canonical BNT hypotheses on both $\mu^A,A$ and $\mu^B,B$,
\leanid{EventuallyNonzeroProportionalMPV}\textsubscript{2} between the
assembled tensors, an index $k_0:\mathrm{Fin}\,r_B$, and the hypothesis
$\forall j,\langle V^{(N)}(A_j),V^{(N)}(B_{k_0})\rangle\to 0$, derive
$\bot$.
The old one-copy fixed-block declarations for this target have been retired.
The corrected \leanid{SectorBNT} proportional matching theorem keeps the raw
coefficients of \leanid{SectorDecomposition} visible and uses the API lemmas in
\doclink{TNLean/MPS/FundamentalTheorem/SectorBNT/Api}.  It does not require a
separate universally quantified fixed-block projection theorem.

\section{Where the per-block projection fails for non-dominant
\texorpdfstring{$k_0$}{k0}}

Projecting \eqref{eq:prop} onto $V^{(N)}(B_{k_0})$ gives, with
$\mathrm{LHS}_N:=\langle V^{(N)}(A_{\mathrm{tot}}),V^{(N)}(B_{k_0})\rangle$
and $\mathrm{RHS}_N:=c_N\langle V^{(N)}(B_{\mathrm{tot}}),V^{(N)}(B_{k_0})\rangle$,
\begin{align}
    \mathrm{LHS}_N
       &=\sum_{j}(\mu^A_j)^N\langle V^{(N)}(A_j),V^{(N)}(B_{k_0})\rangle,
    & \mathrm{RHS}_N
       &=c_N\sum_{k}(\mu^B_k)^N\langle V^{(N)}(B_k),V^{(N)}(B_{k_0})\rangle.
    \label{eq:LHSRHS}
\end{align}
With $|\mu^A_j|\le 1$, every $(\mu^A_j)^N$ is bounded; together with
\eqref{eq:Hk0} this forces $\mathrm{LHS}_N\to 0$. For $\mathrm{RHS}_N$,
the surviving cross-overlap decay statement
\leanid{MPSTensor.IsBNTCanonicalForm.cross_overlap_basis_tendsto_zero}
gives $\langle V^{(N)}(B_k),V^{(N)}(B_{k_0})\rangle\to 0$ for $k\ne k_0$,
while the diagonal block has a nonzero limiting coefficient, so
$\mathrm{RHS}_N=c_N(\mu^B_{k_0})^N(\ell+o(1))$.
In the dominant block, the proportionality identity and the BNT normalizations
give the dominant-adjusted scalar limit
$\bigl|c_N(\mu^B_0/\mu^A_0)^N\bigr|\to 1$.  Therefore
\begin{align}
    \bigl|c_N(\mu^B_{k_0})^N\bigr|
    &\sim
    |\mu^A_0|^N
    \bigl|\mu^B_{k_0}/\mu^B_0\bigr|^N.
    \label{eq:decay}
\end{align}
Under $|\mu^A_0|=1$ and the retired strict ordering, the second factor
in~\eqref{eq:decay} decays geometrically whenever $k_0\ne 0$, so
$\mathrm{RHS}_N\to 0$ together with $\mathrm{LHS}_N$. No contradiction
is available from the per-block projection. The dominant case
$k_0=0$ is the unique index for which the ratio in~\eqref{eq:decay}
is the constant~$1$ and the bound stays uniform; that is the analytic
content that the dominant fixed-block statement captured before the
retirement of the old one-copy layer.

Two independent probes confirm the obstruction. Path~$\alpha$
(\path{docs/audits/2026-05-13_cpsv16_ft_exact_leading_coeff.md},
\ghpr{1654}) shows an exact eventual leading-coefficient identity
$c_N(\mu^B_0)^N\zeta^N=(\mu^A_0)^N$ fails at $r_A=r_B=2$ with distinct
sub-dominant moduli. Path~$\beta$ scope adjustment
(\path{docs/audits/2026-05-13_cpsv16_ft_path_beta_scout.md},
\ghpr{1664}) shows the lifted per-block projection on the
\leanid{SectorDecomposition} surface inherits the same decay
mismatch.

\section{The fixed-block source claim and the formal repair}

The Step~1 passage at
\cite[line~1182]{Cirac2016MPDO_arXiv} asserts the fixed-block contradiction
described above.  Isolating $V^{(N)}(B_{k_0})$ in~\eqref{eq:prop} gives
\begin{align}
    \sum_{j}(\mu^A_j)^N V^{(N)}(A_j)
    -c_N\!\!\!\sum_{k\ne k_0}\!\!(\mu^B_k)^N V^{(N)}(B_k)
    &=c_N(\mu^B_{k_0})^N V^{(N)}(B_{k_0}).
    \label{eq:isolate}
\end{align}
The single-block linear independence of
$\{V^{(N)}(A_j)\}_j\cup\{V^{(N)}(B_{k_0})\}$ does not control the remaining
vectors $V^{(N)}(B_k)$ on the left of~\eqref{eq:isolate}.  Thus it does not by
itself justify coefficient comparison in the full proportionality identity.
The source passage supplies neither all-pairwise cross-family decay nor a
full combined-family linear-independence statement.

The formal development replaces this step rather than formalizing it literally.
For a chosen $P$-block it collects the $P$-blocks whose overlaps with every
$Q$-block decay.  The complementary $P$-blocks already admit scalar relations
with matched $Q$-blocks.  It then applies
\leanid{MPSTensor.restricted_combined_family_eventually_li} to the decaying
$P$-subfamily together with all $Q$-blocks and compares coefficients in that
restricted family.  The selected coefficient is forced to vanish eventually,
contradicting
\leanid{MPSTensor.IsBNTCanonicalForm.coeff_not_eventually_zero}.  The unrestricted
lemma \leanid{MPSTensor.IsBNTCanonicalForm.combined_family_eventually_li} is the
full-family special case of the same repair.  This coefficient-comparison
argument proves the current active-sector matching theorem without using the
retired per-block projection.

\section{Classification}

\emph{Historical proof archaeology and forbidden-route explanation.}  The
retired dominant fixed-block declarations handled only the leading block on the
old one-copy strict-modulus surface.  Extending that projection argument to an
arbitrary block fails because both sides of the projected proportionality may
decay together.  This remains a mathematical warning about the method, not an
unproved theorem required by the current development.

The source line~1182 fixed-block claim and the current formal repair must be
kept distinct.  The repair uses restricted combined-family linear independence,
scalar relations for the complementary blocks, and the raw power sums
$\sum_q\mu_{j,q}^N$ carried by \leanid{SectorDecomposition}.  Together with
\leanid{MPSTensor.IsBNTCanonicalForm.coeff_not_eventually_zero}, this gives the
active-sector coefficient comparison proved in the current theorem.  Collapsing
a sector to one strictly ordered weight destroys precisely the information
needed for that argument, so non-dominant per-block projection is a route that
must not be reintroduced.

\bibliographystyle{alpha}
\bibliography{references}

\end{document}
